\chapter{Suites numériques}
\thispagestyle{fancy}

Dès l'Antiquité, Archimède de Syracuse (-287 ; -212), met en oeuvre une procédure itérative pour trouver une approximation du nombre $\pi$. Il encadre le cercle par des polygones inscrits et circonscrits possédant un nombre de côtés de plus en plus grand. Par ce procédé,
Archimède donne naissance, sans le savoir, à la notion de suite numérique.\\
Vers la fin du XVII$^\text{e}$ siècle, des méthodes semblables sont utilisées pour résoudre des équations de façon approchée pour des problèmes de longueurs, d'aires, ...\\
Un formalisme plus rigoureux de la notion de suite n'apparaitra qu'au début du XIX$^\text{e}$ siècle avec le mathématicien français Augustin Louis Cauchy (1789 ;1857).

\section{Définition et représentation graphique}

\subsection{Définition d'une suite numérique}

\begin{exempleTitre}{Exemple d'introduction}
On considère une liste de nombres formée par tous les nombres impairs rangés dans l'ordre croissant : $1,3,5,7, \ldots$\\
On note $\left(u_{n}\right)$ l'ensemble des éléments de cette suite de nombres tel que :\\
$u_{0}=1, u_{1}=3, u_{2}=5, u_{3}=7, \ldots$\\
On a ainsi défini une suite numérique.\\
On peut lui associer une fonction définie sur $\mathbb{N}$ par 
$$u: \mathbb{N} \rightarrow \mathbb{R}$$
$$\phantom{......................}n \mapsto u(n)=u_{n}$$

\end{exempleTitre}

\begin{definitions}
\begin{itemDef}
\item Une \grasInDef{suite numérique} $\left(u_{n}\right)$ est une liste ordonnée de nombres réels telle qu'à tout entier $n$ on associe un nombre réel noté $u_{n}$.
\item $u_{n}$ est appelé \grasInDef{le terme de rang $\boldsymbol{n}$} de cette suite (ou d'indice $\left.n\right)$.
\end{itemDef}
\end{definitions}

\subsection{Suite définie par une formule explicite}

\begin{exemples}
\begin{itemEx}
\item Pour tout $n$ de $\mathbb{N}$, on donne : $u_{n}=2 n$ qui définit la suite des nombres pairs. Les premiers termes de cette suite sont donc :\\
$u_{0}=2 \times 0=0$\\
$u_{1}=2 \times 1=2$\\
$u_{2}=2 \times 2=4$\\
$u_{3}=2 \times 3=6$
\item Pour tout $n$ de $\mathbb{N}$, on donne $: v_{n}=3 n^{2}-1$ Les premiers termes de cette suite sont donc :\\
$v_{0}=3 \times 0^{2}-1=-1$\\
$v_{l}=3 \times 1^{2}-1=2$\\
$v_{2}=3 \times 2^{2}-1=11$\\
$v_{3}=3 \times 3^{2}-1=26$
\end{itemEx}
\end{exemples}
\begin{definition}
    Lorsqu'on génère une suite par une formule explicite, chaque terme de la suite est exprimé en fonction de $n$ et indépendamment des termes précédents.
\end{definition}

\subsection{Suite définie par une relation de récurrence}
\begin{exemples}
\begin{itemEx}
\item On définit la suite $\left(u_{n}\right)$ par: $u_{0}=5$ et chaque terme de la suite est le triple de son précédent. Les premiers termes de cette suite sont donc :\\
$u_{0}=5$\\
$u_{1}=3 \times u_{0}=3 \times 5=15$\\
$u_{2}=3 \times u_{1}=3 \times 15=45$\\
De façon générale, on peut noter $: u_{n+1}=3 u_{n}$
\item On définit la suite $\left(v_{n}\right)$ par : $v_{0}=3$ et pour tout $n$ de $\mathbb{N}, v_{n+1}=4 v_{n}-6$ Les premiers termes de cette suite sont donc :\\
$v_{0}=3$\\
$v_{1}=4 v_{0}-6=4 \times 3-6=6$\\
$v_{2}=4 v_{1}-6=4 \times 6-6=18$\\
$v_{3}=4 v_{2}-6=4 \times 18-6=66$\\
Contrairement à une suite définie par une formule explicite, il n'est pas possible, dans l'état, de calculer par exemple $v_{13}$ sans connaître $v_{12}$.\\
\ \\
Cependant il est possible d'écrire un algorithme avec Python :
\begin{center}
    \includegraphics[scale=0.4]{1.08.00-Img01}
\end{center}
Ou sur une calculatrice : 
\begin{center}
    \includegraphics[scale=0.4]{1.08.00-Img02}
\end{center}
\end{itemEx}
\end{exemples}
\begin{definition}
Lorsqu'on génère une suite par une relation de récurrence, chaque terme de la suite s'obtient à partir de son terme précédent.Il faut connaître au moins un terme de la suite pour calculer les autres termes.
\end{definition}

\begin{remarque}
Le mot \emph{récurrence} vient du latin \emph{recurrere} qui signifie "revenir en arrière".
\end{remarque}

\subsection{Représentation graphique d'une suite}

Dans un repère du plan, on représente une suite par un nuage de points de coordonnées $\left(n ; u_{n}\right) .$
\begin{exemple}
Pour tout $n$ de $\mathbb{N}$, on donne $: u_{n}=\dfrac{n^{2}}{2}-3$.\\
On construit le tableau de valeurs avec les premiers termes de la suite :
\begin{tabular}{|Xc|Xc|Xc|Xc|Xc|Xc|Xc|Xc|Xc|Xc|}
\hline$n$ & 0 & 1 & 2 & 3 & 4 & 5 & 6 & 7 & 8 \\
\hline$u_{n}$ & $-3$ & $-2,5$ & $-1$ & 1,5 & 5 & 9,5 & 15 & 21,5 & 29 \\
\hline
\end{tabular}\\
Il est possible d'obtenir un nuage de points à l'aide d'un logiciel.
\begin{center}
    \includegraphics[scale=0.5]{1.08.00-Img03}
\end{center}

\end{exemple}

\section{Sens de variation d'une suite numérique}

\begin{exemple}
On a représenté ci-dessous le nuage de points des premiers termes d'une suite $\left(u_{n}\right)$ :
\begin{center}
    \includegraphics[scale=0.5]{1.08.00-Img04}
\end{center}
On peut conjecturer que cette suite est croissante. On constate par exemple que $u_{1}<u_{2}$ ou encore $u_{4}<u_{5}$. De manière générale, on peut écrire : $u_{n}<u_{n+1}$
\end{exemple}

\begin{definitions}
Soit une suite numérique $\left(u_{n}\right)$.
\begin{itemDef}
\item La \grasInDef{suite $\left(u_{n}\right)$ est croissante} signifie que pour tout entier $n$, on a $u_{n+1} \geqslant u_{n}$.
\item La \grasInDef{suite $\left(u_{n}\right)$ est décroissante} signifie que pour tout entier $n$, on a $u_{n+1} \leqslant u_{n}$.
\end{itemDef}
\end{definitions}

\begin{methodeTitre}{Étudier les variations d'une suite}
\grasInMet{Énoncé:}
\begin{enumMet}
\item Pour tout $n$ de $\mathbb{N}$, on donne la suite $\left(u_{n}\right)$ définie par : $u_{n+1}=u_{n}+2$.\\
Démontrer que la suite $\left(u_{n}\right)$ est croissante.
\item Pour tout $n$ de $\mathbb{N}$, on donne la suite $\left(v_{n}\right)$ définie par : $v_{n}=4 n+4$.\\
Démontrer que la suite $\left(v_{n}\right)$ est croissante.
\end{enumMet}
\grasInMet{Réponse:}
\begin{enumMet}
\item $u_{n+1}-u_{n}=2>0$\\
On en déduit que $\left(u_{n}\right)$ est croissante.
\item On commence par calculer la différence $v_{n+1}-v_{n}$ :\\
On a $: v_{n}=4 n+4$ donc $v_{\mVert{n+1}}=4\mVert{(n+1)}+4=4 n+4+4=4 n+8$
\begin{align*}
v_{n+1}-v_{n}&=4 n+8-(4 n+4) \\
&=4 n+8-4 n-4 \\
&=4
\end{align*}
On étudie ensuite le signe de $v_{n+1}-v_{n}$ :\\
Pour tout $n$ entier $v_{n+1}-v_{n}=4>0$.\\
On en déduit que la suite $\left(v_{n}\right)$ est croissante.
\end{enumMet}
\end{methodeTitre}

